\documentclass[journal, twocolumn]{IEEEtran}
\usepackage[UTF8]{ctex}
\usepackage{amsmath, amssymb}
\usepackage{graphicx}
\usepackage{booktabs}
\usepackage{multirow}
\usepackage{listings}
\usepackage{xcolor}
\usepackage{geometry}

\lstset{
    basicstyle=\footnotesize\ttfamily,
    breaklines=true,
    breakatwhitespace=false,
    language=C++,
    keywordstyle=\color{blue},
    commentstyle=\color{gray},
    numbers=left,
    numberstyle=\tiny\color{gray},
    frame=single,
    aboveskip=3pt,
    belowskip=3pt,
}

\title{\textbf{基于多层次超图划分与改进 FM 算法的 VLSI 网表二分问题求解}}

\author{\textbf{组号：6} \quad 陈奕航 23336047 \quad 陈泽义 23336050 \quad 陈俊鑫 23336040}

\begin{document}

\maketitle

\begin{abstract}
VLSI 网表二分划分是物理设计中的核心问题，要求在满足 2\% 面积平衡约束的前提下最小化割边数量。本文基于多层次超图划分框架，设计并实现了一种结合改进 Fiduccia-Mattheyses (FM) 算法的求解器。算法采用多次重启策略，在粗化阶段引入小网加权与度数惩罚的匹配评分机制，在细化阶段使用基于最大堆和延迟更新的增量式 FM 变体，并针对不同层次配置差异化参数与打磨过程。实验在 18 个 IBM ISPD 基准测试用例上进行，结果在中小规模实例上接近最优，验证了方案的有效性。
\end{abstract}

\section{问题描述}

超图划分是 VLSI CAD 流程中的基础问题。在芯片物理设计中，网表的二分划分质量直接影响布局布线的效果。形式化地，给定超图 $G = (V, E)$，其中 $V$ 为节点（逻辑单元）集合，$E$ 为超边（线网）集合，目标是将 $V$ 划分为两个子集 $X$ 和 $Y$，使得被切割的超边数最小，同时满足：

\begin{equation}
\frac{1}{2} - \epsilon \le \frac{|X|}{|V|} \le \frac{1}{2} + \epsilon, \quad \epsilon = 2\%
\end{equation}

该问题是 NP-Hard 的，实际求解依赖启发式算法。最经典的 FM 算法通过贪心移动节点来迭代优化割边数，但单层 FM 在大规模图上效果有限。为此，本文采用多层次（multilevel）框架——先通过节点匹配合并逐步缩减图规模（粗化），在粗图上进行初始划分并逐层细化回原图，每层均辅以 FM 局部优化。

\section{数据结构设计}
\label{sec:datastruct}

\subsection{原始图表示}

实验框架提供了基础的面向对象图表示，包含三个核心类：

\textbf{Node} 类记录每个节点的唯一索引 \texttt{index} 及其关联的所有线网列表 \texttt{vector<Net*> nets}。\textbf{Net} 类记录线网索引及该线网连接的所有节点列表 \texttt{vector<Node*> nodes}。\textbf{Graph} 类统一管理所有节点和线网，提供 \texttt{get\_or\_create\_node} 和 \texttt{add\_net} 接口，内部使用哈希表 \texttt{unordered\_map} 实现索引到对象的快速映射。这种设计使得遍历一个节点的所有邻居（通过节点 $\to$ 线网 $\to$ 节点的两跳访问）的复杂度与度数成正比。

\subsection{超图内部表示}

为提升算法效率，我们设计了紧凑的 \texttt{Hypergraph} 结构体作为内部计算表示：

\begin{lstlisting}
struct Hypergraph {
    int n;                          // 节点总数
    vector<vector<int>> nets;       // nets[i]: 线网i上的所有节点
    vector<vector<int>> node_nets;  // node_nets[u]: 节点u关联的所有线网
    vector<int> node_weight;        // 节点权重（粗化时累加）
};
\end{lstlisting}

该结构的核心设计思想是\textbf{双向索引}——同时维护“线网到节点”和“节点到线网”两组邻接关系。这使得以下操作均为 $O(d)$ 复杂度（$d$ 为节点度数）：

\begin{itemize}
    \item \textbf{访问节点的所有邻居}：遍历 \texttt{node\_nets[u]} 找到所有关联线网，再遍历每个线网的 \texttt{nets[net\_idx]} 获得同网所有节点。
    \item \textbf{计算节点移动的割边增量}：遍历 \texttt{node\_nets[u]}，对每条关联线网检查其两侧节点计数的变化，$O(d)$ 即可完成。
    \item \textbf{粗化匹配评分}：遍历 \texttt{node\_nets[u]} 获得所有候选邻居，高效累加匹配得分。
\end{itemize}

在 FM 算法的增量更新阶段，当一个节点被移动后，仅需通过 \texttt{node\_nets[u]} $\to$ \texttt{nets[net\_idx]} 路径定位受影响邻居，利用时间戳去重后更新增益值。这种双向索引设计避免了全图扫描，是算法性能的关键保障。

此外，\texttt{node\_weight} 向量在粗化过程中记录每个粗化节点所代表的原始节点数量，用于平衡约束检查。原始图中每个节点权重为 1，粗化后的权重为组内成员权重之和。

\section{多层次划分框架}

\subsection{总体流程}

如图 \ref{fig:flow} 所示，算法采用经典的 V-cycle 多层次框架，每次重启（共 8 次）执行一组完整的粗化-初始划分-细化流程。

\begin{figure}[htbp]
\centering
\fbox{\begin{minipage}{0.85\columnwidth}
\footnotesize
\centering
\textbf{多次重启}（8次）\\[2pt]
\fbox{\begin{minipage}{0.9\columnwidth}
\centering
粗化循环（最多5层）\\
$\downarrow$\\
最粗层初始划分（3种策略竞争）\\
$\downarrow$\\
FM 细化 + 打磨\\
$\downarrow$\\
逐层反粗化 + 分层 FM\\
$\downarrow$\\
保留最优解
\end{minipage}}
\end{minipage}}
\caption{总体算法流程}
\label{fig:flow}
\end{figure}

粗化阶段通过节点匹配逐步缩减图规模，将原始超图压缩至粗化层次。当粗化到最粗层后，在粗图上执行初始划分并优化，然后通过反粗化将划分逐层投影回更细的图，每层均运行 FM 算法精化结果。最终在最细层（原始图）得到最终划分并计算割边数。

\subsection{粗化机制：基于评分的节点匹配}

粗化是决定最终解质量的关键。每次粗化调用 \texttt{coarsen\_once} 函数，将当前层节点配对合并。

\subsubsection{节点访问顺序}

首先将所有节点随机打乱（\texttt{shuffle}），然后按度数降序稳定排序。高度数节点优先匹配，因为这些节点连接线网多，对割边的影响更大，优先确定其归属有助于减少后续优化的不确定性。

\subsubsection{匹配评分函数}

对于未匹配节点 $u$，遍历其所有关联线网。对每条线网的每个未匹配节点 $v$，累加匹配得分：

\begin{equation}
\text{score}(u, v) = \sum_{\text{net} \in \text{nets}(u) \cap \text{nets}(v)} \frac{w(\text{net})}{\text{deg\_penalty}(v)}
\end{equation}

其中各项的定义体现了三个关键设计考量：

\textbf{线网权重 $w(\text{net})$}：定义为 $\frac{1}{|\text{net}| - 1}$。线网越小（连接节点越少），权重越高，因为小线网代表的耦合关系更强。例如，连接 2 个节点的线网权重为 1.0，而连接 64 个节点的线网权重仅为 $1/63 \approx 0.016$。

\textbf{小网奖励}：当线网大小 $\le 4$ 时，权重额外乘以系数 $\texttt{kSmallNetBoost} = 1.7$。这种设计基于观察：在真实 VLSI 网表中，小线网（2-4 个节点）是最常见的线网类型，它们代表了设计中紧密的局部连接，应当被优先保留在同一划分侧。

\textbf{度数惩罚}：$\text{deg\_penalty}(v) = 1.0 + \texttt{kMatchDegreePenalty} \cdot |\text{nets}(v)|$，其中 $\texttt{kMatchDegreePenalty} = 0.06$。该惩罚项抑制了与高度数节点（“超级枢纽”）的匹配倾向。若不加以惩罚，高度数节点会因其大量连接而吸收过多匹配，形成权重极高的粗化节点，导致粗化图丧失结构信息。适度惩罚后，匹配更倾向于度数相近的节点对。

\textbf{同分决胜}：得分相同时，选择度数更小的节点作为匹配对象，避免形成过大的超级节点，保持粗化图的结构平衡。

\subsubsection{粗化图构建}

得分最高的节点对 $\{u, v_{\text{best}}\}$ 形成一对（若无可配对节点则 $u$ 单独成组）。每组生成一个粗化节点：
\begin{itemize}
    \item 粗化节点权重 = 组内成员权重之和
    \item 粗化线网：将原线网中节点替换为其对应的粗化节点索引，去重后若仍含 $\ge 2$ 个节点则保留
\end{itemize}

粗化停止条件包括：达到最大层数 5、节点数降至 320 以下、粗化率 $> 0.92$（缩减不充分）。

\subsection{初始划分的三种策略}

在最粗层上，算法采用三种不同策略生成初始划分，每种均需满足 $[0.48, 0.52]$ 的权重平衡约束，最终选取 FM 优化后割边最少的结果。

\subsubsection{基于节点度的平衡划分}

按度数降序排列节点，贪心地分配到 X 或 Y 侧。高度数节点（连接线网多）优先放置，因为它们在后续细化中有更大的优化操作空间。优先将节点放入当前权重较小的一侧以维持平衡。该策略的优点是计算简单高效（仅需排序一次），且度数高的节点对割边影响大，优先分配使初始解已有一定质量。

\subsubsection{随机平衡划分}

完全随机打乱节点顺序，然后将每个节点贪心放入较轻的一侧（权重相等时抛硬币）。此策略不依赖于任何结构信息，旨在提供\textbf{多样性}。由于后续 FM 可能陷入局部最优，随机初始化提供了不同的搜索起点，增大了覆盖全局最优的概率。

\subsubsection{基于小网权重的平衡划分}

为每个节点计算一个小网亲和度分数：

\begin{equation}
\text{score}(u) = \sum_{\text{net} \in \text{nets}(u)} \frac{1}{|\text{net}| - 1} \cdot \mathbf{1}_{|\text{net}| \le 4} \cdot \texttt{kSmallNetBoost}
\end{equation}

按分数降序排列后贪心分配，同时追踪两侧的累计分数（$s_x$ 和 $s_y$），在权重允许时优先将节点放入累计分数较小的一侧，使小网相关节点分布更均匀。该策略的优点是直接以割边相关的度量指标指导划分，与后续 FM 优化目标一致。

三者的关系可总结为：度数策略提供基于结构信息的确定性起点，随机策略提供多样性探索，小网策略提供与优化目标高度对齐的贪心起点。在 8 次重启中，每种策略各被尝试 $8 \times 1$ 次，结果择优保留。

\section{改进的 FM 细化算法}

\subsection{经典 FM 算法回顾}

经典 FM 算法通过重复“选择增益最大的可移动节点 $\to$ 移动并锁定 $\to$ 回退到最佳前缀”的单轮迭代，逐步减小割边数。其时间复杂度在桶结构（bucket）实现下可达每轮 $O(P)$，$P$ 为总引脚数。

\subsection{基于最大堆与延迟更新的变体设计}

本实现中使用最大堆（\texttt{priority\_queue}）替代传统桶数组，并结合延迟更新（lazy update）机制：

\begin{lstlisting}
struct GainEntry {
    int gain;     // 移动收益
    int node;     // 节点索引
    int version;  // 版本号
};
struct GainEntryCmp {
    bool operator()(const GainEntry &a, const GainEntry &b) const {
        return a.gain < b.gain;  // 最大堆
    }
};
\end{lstlisting}

\subsection{核心操作流程}

\textbf{1) 初始化增益}：计算每个节点的初始增益 $g(u) = -\Delta_{\text{cut}}(u)$（$\Delta_{\text{cut}}$ 为移动导致的割边变化量，负值表示改善），将所有节点推入堆中。

\textbf{2) 增量更新}：当一个节点 $u$ 被移动后，仅重新计算其“受影响邻居”的增益。受影响邻居定义为通过至少一条线网与 $u$ 相连且未被锁定的节点。使用时间戳数组 \texttt{dedup\_mark} 确保同一轮中每个邻居仅被更新一次，避免因多条共享线网导致重复计算。受影响的节点获得新的增益值、版本号递增、并被重新推入堆中。

\textbf{3) 延迟更新}：从堆顶取出条目时，比较 \texttt{entry.version} 与当前 \texttt{gain\_version[node]}。若版本不匹配，说明该条目的增益已过期，直接丢弃并继续取下一个。这避免了每次移动后重建整个堆的开销。

\textbf{4) 平衡约束过滤}：取出有效条目后，检查移动该节点是否会破坏平衡约束 $[W_{\min}, W_{\max}]$。若不满足，暂存到延迟数组（\texttt{deferred}），待选出合法候选后再将这些条目重新推回堆中，避免丢失有效候选。

\textbf{5) 全量扫描兜底}：若堆取样范围内（受 \texttt{candidate\_size} 限制）无合法候选，退化执行全图 $O(n)$ 扫描，确保不会因堆顶积压过期条目而遗漏可用的移动。

\textbf{6) 前缀最佳回滚}：移动序列完成后，找到累计增益达到最大值的移动前缀。若最佳前缀增益 $\le 0$（本轮无实质性改善），回退全部移动；否则仅保留前缀部分，回退后续移动。

\subsection{增量增益计算}

\texttt{delta\_cut\_for\_move} 函数计算单个节点移动对割边数的影响。对节点的每条关联线网，比较移动前后的“是否被割”状态：

\begin{itemize}
    \item 移动前：$\text{before\_cut} = (\text{count\_x} > 0 \land \text{count\_y} > 0)$
    \item 移动后：按移动方向更新 $\text{count\_x}$ 和 $\text{count\_y}$，重新判断
    \item $\Delta_{\text{cut}} = \sum (\text{after\_cut} - \text{before\_cut})$
\end{itemize}

每条线网的贡献仅为 $-1, 0, +1$ 三者之一，该计算为 $O(\text{degree}(u))$ 复杂度。

\subsection{分层参数配置}

针对不同粗化层的特点，算法配置了差异化的 FM 参数，如表 \ref{tab:fmlayers} 所示。

\begin{table}[htbp]
\centering
\caption{FM 分层参数配置}
\label{tab:fmlayers}
\footnotesize
\begin{tabular}{lccc}
\toprule
\textbf{参数} & \textbf{粗层} & \textbf{中层} & \textbf{细层} \\
\midrule
FM 轮数 & 8 & 8 & 10 \\
每轮最大步数 & 9000 & 10000 & 14000 \\
堆检查上限 & 1800 & 2200 & 3000 \\
打磨轮数 & 2 & 2 & 3 \\
打磨最大步数 & 2500 & 3000 & 4500 \\
打磨堆检查 & 2500 & 3200 & 4500 \\
\bottomrule
\end{tabular}
\end{table}

\textbf{分层的设计逻辑}：

\textbf{粗层}（最粗化层）：节点数少（数百量级），FM 每轮收敛快。参数设置较为保守，快速跳过优化空间较小的粗化层。

\textbf{中层}（中间粗化层）：节点数中等（数千量级），适度增大步数和堆检查范围，在计算量可控的前提下提供更充分的局部搜索。

\textbf{细层}（原始图，即 level 0）：节点数最大（可达 21 万），是决定最终割边质量的关键层。因此配置最激进的参数——更多的 FM 轮数、更大的步数和堆检查范围，充分挖掘优化潜力。

\textbf{打磨过程（Polish Passes）}：在每层主 FM 优化之后，额外追加若干轮打磨 FM，采用更大的候选搜索范围和更多的步数限制。打磨轮的定位是“在已有较好解的基础上，用更大搜索范围试探是否存在被主轮遗漏的增益机会”，能够在时间预算允许时进一步压缩割边数。

\subsection{与其他 FM 变体的比较}

相较于经典 bucket-FM（使用双向桶数组按增益索引节点），堆变体有以下权衡：

\textbf{优势}：
\begin{itemize}
    \item 支持可变节点权重，无需预知增益范围
    \item C++ STL 的 \texttt{priority\_queue} 高度优化，常数因子小
    \item 延迟更新避免了全堆重建，增量更新仅影响局部节点
\end{itemize}

\textbf{劣势}：
\begin{itemize}
    \item 每次 pop 为 $O(\log n)$，而 bucket 的查找为 $O(1)$
    \item 过期条目累积会增加无效 pop 次数，需设置堆检查上限
\end{itemize}

在大规模实例上，堆方式的 $O(\log n)$ 开销可通过限制堆检查次数和增量更新范围来控制，整体性能可接受。

\section{外层控制策略}

\subsection{多次重启}

设置 $\texttt{kRestarts} = 8$ 次完整重启。每次重启使用不同的随机种子（通过 \texttt{mt19937} 伪随机数生成器控制），产生不同的粗化匹配和初始划分序列。多次重启的核心价值在于：

\begin{itemize}
    \item 粗化阶段中的随机打乱导致每次的节点匹配顺序不同，粗化图的结构因而存在差异，为细化阶段提供了不同的搜索起点
    \item 初始划分中的随机策略在每次重启中产生不同的起点
    \item 8 次重启在质量和耗时之间取得了平衡——实验表明继续增加重启次数收益递减
\end{itemize}

\subsection{严格时间控制}

算法设置总时间预算 $\texttt{kTotalTimeBudgetSeconds} = 90$ 秒。通过 \texttt{chrono::steady\_clock} 在以下关键节点检查时间：

\begin{itemize}
    \item 每次重启迭代开始前
    \item 每次粗化层构建前
    \item 每次初始划分策略切换前
    \item 每轮 FM 及每步节点移动中
\end{itemize}

一旦超时，立即中断当前路径，返回已保存的最优解。这种细粒度的时间检查确保了即使在大规模实例上算法也能在 90 秒内完成，同时保留在此之前找到的最佳结果。

\section{实验结果分析}

\subsection{测试环境与配置}

实验在 18 个 IBM ISPD 基准测试用例上进行，节点规模从 12,752 到 210,613 不等。平衡容差 $\epsilon = 2\%$。参考最优解取自带数据集中 $\epsilon=2\%$ 列的理论下界。

\subsection{整体结果}

表 \ref{tab:results} 汇总了各测试用例的割边结果及与最优解的对比。

\begin{table}[htbp]
\centering
\caption{实验结果汇总}
\label{tab:results}
\footnotesize
\begin{tabular}{lrrrrr}
\toprule
\textbf{用例} & \textbf{节点数} & \textbf{本文结果} & \textbf{最优(2\%)} & \textbf{差值} & \textbf{相对差} \\
\midrule
IBM01 & 12,752 & 261 & 200 & 61 & 30.50\% \\
IBM02 & 19,601 & 355 & 307 & 48 & 15.64\% \\
IBM03 & 23,136 & 1074 & 951 & 123 & 12.93\% \\
IBM04 & 27,507 & 639 & 573 & 66 & 11.52\% \\
IBM05 & 29,347 & 1848 & 1706 & 142 & 8.32\% \\
IBM06 & 32,498 & 1074 & 962 & 112 & 11.64\% \\
IBM07 & 45,926 & 1044 & 878 & 166 & 18.91\% \\
IBM08 & 51,309 & 1219 & 1140 & 79 & 6.93\% \\
IBM09 & 53,395 & 1035 & 620 & 415 & 66.94\% \\
IBM10 & 69,429 & 1582 & 1253 & 329 & 26.26\% \\
IBM11 & 70,558 & 1368 & 1051 & 317 & 30.16\% \\
IBM12 & 71,076 & 2066 & 1919 & 147 & 7.66\% \\
IBM13 & 84,199 & 967 & 831 & 136 & 16.37\% \\
IBM14 & 147,605 & 3468 & 1842 & 1626 & 88.27\% \\
IBM15 & 161,570 & 5595 & 2730 & 2865 & 104.95\% \\
IBM16 & 183,484 & 5705 & 1827 & 3878 & 212.26\% \\
IBM17 & 185,495 & 3321 & 2270 & 1051 & 46.30\% \\
IBM18 & 210,613 & 5479 & 1521 & 3958 & 260.22\% \\
\midrule
\textbf{平均} & -- & -- & -- & -- & \textbf{54.21\%} \\
\bottomrule
\end{tabular}
\end{table}

\subsection{结果分析}

\textbf{中小规模表现良好}：在 IBM02-IBM08、IBM12-IBM13 等中小规模实例上，相对差距控制在 7\%--19\% 之间，IBM08（6.93\%）和 IBM12（7.66\%）表现尤为突出。这些实例的节点数在 2--8 万之间，5 层粗化足以充分压缩图规模，FM 细化有足够的时间完成多轮迭代。这验证了多层次框架在中等规模问题上的有效性。

\textbf{大规模实例退化明显}：IBM14-IBM18（节点数 14--21 万）上差距急剧扩大至 88\%--260\%。主要原因有三：

\begin{enumerate}
    \item \textbf{粗化质量瓶颈}：大规模图上，基于局部评分的贪心匹配难以捕获全局结构信息，粗化图可能丢失关键的超边分布特征，导致初始划分的搜索空间偏离最优区域。
    \item \textbf{细化深度不足}：90 秒时间预算下，大规模图的最细层 FM 仅能完成有限轮迭代，增量更新的局部搜索半径不足以逃离较深的局部最优。
    \item \textbf{重启策略的局限}：8 次重启在 90 秒内，大规模图每次重启的粗化和细化都不充分，不同随机种子产生的解多样性有限。
\end{enumerate}

\textbf{IBM09 的异常}：IBM09（53,395 节点，相对差 66.94\%）是中等规模中表现最差的。其线网数（60,902）相对于节点数偏高，表明其超图结构更密集。密集图中的节点移动增益计算更加复杂，局部搜索更容易陷入振荡。

\textbf{IBM05 的突出表现}：相对差仅 8.32\%，是 2 万节点以上实例中最好的成绩。该实例的线网数（28,446）略少于节点数（29,347），结构相对稀疏，割边数较大（1848 vs 1706），说明在稀疏图上本算法能接近最优。

\subsection{参数敏感性讨论}

通过实验观察，以下参数对结果影响较大：

\begin{itemize}
    \item \textbf{小网奖励系数}（1.7）：粗化匹配中对小线网的保护程度。系数过大会导致匹配过于保守（仅配对直接相连节点），过小则弱化局部耦合信息。
    \item \textbf{度数惩罚系数}（0.06）：控制超级节点形成的抑制程度。惩罚过小会导致粗化图出现少数超大权重节点，使后续平衡约束更难满足。
    \item \textbf{粗化层数}：5 层对中小实例足够，但对大实例可考虑加深至 7-8 层以获得更充分的规模缩减。
    \item \textbf{FM 堆检查上限}：在效率和搜索完整性之间的折中。3000 的细层上限在多数情况下足够，但大图上可能需要动态调整。
\end{itemize}

\section{改进方向}

\subsection{近期改进}

\textbf{标准 FM Bucket 结构}：将当前的最大堆替换为双向桶（gain bucket）结构，利用增益值的有限范围（$-\max\_\text{degree} \sim +\max\_\text{degree}$）实现 $O(1)$ 的最大增益查找。这将消除堆 $O(\log n)$ 的 pop 开销，大幅提升大规模图上的 FM 吞吐量。结合更高效的可行性过滤，可在相同时间预算内完成更多轮细化。

\textbf{更深的粗化层次}：对大规模实例适当增加粗化层数（如 7-8 层），使粗化图充分缩小，在最粗层上 FM 的搜索空间更紧凑。

\subsection{中期改进}

\textbf{自适应参数策略}：当前所有实例共享同一套参数。未来可根据输入图的规模（节点数、线网数、平均度数）动态调整重启次数、粗化层数、FM 轮数和堆检查范围。例如，对 20 万节点以上的实例自动分配更多粗化层但减少单次重启的 FM 轮数以换取重启多样性。

\textbf{结构感知匹配}：在粗化匹配评分中引入超边分布熵或局部割敏感度度量，使匹配决策不仅基于局部线网大小，还考虑全局结构，减少粗化过程中的信息损失。

\subsection{长期改进}

\textbf{并行粗化与细化}：粗化阶段的节点匹配和 FM 的增益更新均可利用多线程并行加速，特别是大规模图上增量更新的邻居遍历适合并行化。

\textbf{机器学习引导的划分}：利用图神经网络或强化学习预测节点归属倾向或线网切割概率，将其作为初始划分或 FM 增益计算的先验，提升解的初始质量和搜索方向。

\section{总结}

本文实现了一个基于多层次超图划分与改进 FM 算法的 VLSI 网表二分求解器。在数据结构方面，采用双向索引的超图表示实现了邻接关系的高效访问。在算法方面，通过带奖励/惩罚的匹配评分机制保证粗化质量，三种差异化初始划分策略提供搜索多样性，基于最大堆与延迟更新的增量式 FM 变体在效率与质量之间取得平衡，分层参数配置使不同粗化层获得适配的优化力度，多次重启与严格时间控制确保了解的全局优化潜力与运行稳定性。在 18 个 IBM ISPD 基准测试上，中小规模实例的结果接近最优，验证了方案的有效性；大规模实例上的性能差距也为后续改进指明了方向。

% \bibliographystyle{IEEEtran}
% \nocite{*}
% \bibliography{refs}

\end{document}
